\begin{solution}
Let $s,t$ be two gluings. For every $i$, $s|_{U_i}=s_i=t|_{U_i}$. Locality gives $s=t$.
\end{solution}

\begin{solution}
For the empty cover of $\varnothing$, there is exactly one compatible family. The sheaf axiom says there is exactly one section over $\varnothing$. In abelian groups the unique one-element group is $0$; in sets it is a singleton.
\end{solution}

\begin{solution}
Given compatible $f_i\in C^0(U_i,\RR)$, define $f(x)=f_i(x)$ for $x\in U_i$. Compatibility makes this well defined. Each $f|_{U_i}=f_i$ is continuous, so $f$ is continuous. Pointwise equality gives locality.
\end{solution}

\begin{solution}
Compatible smooth functions glue pointwise. Around every point the glued function equals one of the smooth local functions, hence is smooth. Locality is pointwise equality.
\end{solution}

\begin{solution}
Compatible locally constant maps glue pointwise. Around $x$, choose a cover member and a neighborhood on which its local section is constant; the glued section equals it there, hence is locally constant.
\end{solution}

\begin{solution}
Take $U=U_1\sqcup U_2$ and distinct $a,b\in A$. The local sections $a$ and $b$ are compatible on the empty overlap, but a global element of the naive constant presheaf is a single $c\in A$, whose restrictions would both equal $c$.
\end{solution}

\begin{solution}
Restrictions preserve boundedness, so this is a presheaf. On $U_n=(-n,n)$ let $f_n(x)=x$. They agree on overlaps and are bounded on each $U_n$. Their unique pointwise gluing is $f(x)=x$, unbounded on $\RR$, so gluing fails.
\end{solution}

\begin{solution}
Reflexivity uses $W=U$. Symmetry is immediate. If $(U,s)\sim(V,t)$ on $W_1$ and $(V,t)\sim(Z,r)$ on $W_2$, then $W_1\cap W_2$ is a neighborhood of $x$ on which $s=t=r$, proving transitivity.
\end{solution}

\begin{solution}
For $f(t)=t$ and $g(t)=0$, $f(0)=g(0)=0$, but no neighborhood of $0$ has $f=g$, so $f_0\neq g_0$.
\end{solution}

\begin{solution}
Take $f=0$ and a smooth bump $g$ supported in $(2,3)$. Then $g$ vanishes on a neighborhood of $0$, so $f_0=g_0$, while $g\neq0$ globally.
\end{solution}

\begin{solution}
For germs $s_x,t_x$, choose representatives on $U,V$ and define sum/product using restrictions to $U\cap V$. Changing representatives changes nothing after shrinking to a common neighborhood; hence the operations are well defined and inherit the ring axioms.
\end{solution}

\begin{solution}
A locally constant section has some neighborhood of $x$ on which it equals a constant $a\in A$, so its germ maps to $a$. This is well defined and bijective, with inverse represented by the constant section $a$ on any neighborhood.
\end{solution}

\begin{solution}
At $x$, all neighborhoods receive value $A$ with identity transition maps, so the colimit is $A$. If $y\neq x$, choose an open neighborhood $V$ of $y$ not containing $x$; all smaller neighborhoods map through $0$, so the stalk is $0$.
\end{solution}

\begin{solution}
Suppose $s_x=t_x$. Then on some neighborhood $W$ of $x$, $s|_W=t|_W$. Naturality gives $\varphi(s)|_W=\varphi(t)|_W$, so the image germs agree. Thus the formula is independent of representative.
\end{solution}

\begin{solution}
For a representative $s_x$, $(\psi\circ\varphi)_x(s_x)=(\psi(\varphi(s)))_x=\psi_x(\varphi_x(s_x))$. The identity statement is identical.
\end{solution}

\begin{solution}
If all germs agree, then for every $x$ there is $V_x$ on which $s=t$. The $V_x$ cover $U$, so locality gives $s=t$. The forward implication is immediate.
\end{solution}

\begin{solution}
Apply the equality-from-germs result to $s$ and the zero section. Their germs agree exactly when $s_x=0$ for all $x$.
\end{solution}

\begin{solution}
For $s\in\mathcal F(U)$, the germs of $\varphi_U(s)$ and $\psi_U(s)$ are $\varphi_x(s_x)$ and $\psi_x(s_x)$, hence equal. Sections of $\mathcal G$ are determined by germs, so $\varphi_U(s)=\psi_U(s)$ for all $U,s$.
\end{solution}

\begin{solution}
The singleton cover $U=\bigcup_{x\in U}\{x\}$ has empty pairwise overlaps for distinct points. Any family of singleton sections is compatible and glues uniquely. Since $\mathcal F_x\cong\mathcal F(\{x\})$, the result follows.
\end{solution}

\begin{solution}
Choose groups $A=\mathcal F_p$ and $B=\mathcal F_q$. Then $\mathcal F(\{p,q\})\cong A\times B$, and restrictions are the projections. Conversely any pair $(A,B)$ defines such a sheaf. Morphisms correspond to pairs of group homomorphisms.
\end{solution}

\begin{solution}
A sheaf is data $A=\mathcal F(X)$, $B=\mathcal F(\{\eta\})$, and a restriction map $A\to B$. Since $x$ has no smaller neighborhood than $X$, $\mathcal F_x\cong A$. The neighborhoods of $\eta$ include $\{\eta\}$, terminal in the directed system, so $\mathcal F_\eta\cong B$.
\end{solution}

\begin{solution}
If $s_x=0$, equality of germs means there is a neighborhood $V_x$ on which $s|_{V_x}=0$. Thus the zero-germ locus is the union of such neighborhoods and is open. Its complement is the nonzero-germ locus; the support is its closure.
\end{solution}

\begin{solution}
Every neighborhood of $x$ in $U$ is $U\cap V$ for a neighborhood $V$ in $X$, and these are cofinal in the neighborhood system at $x$. The direct limits defining the two stalks are therefore canonically isomorphic.
\end{solution}

\begin{solution}
Products of compatible families glue componentwise, so $\mathcal F\times\mathcal G$ is a sheaf. A germ of a pair is exactly a pair of germs, giving the canonical stalk isomorphism.
\end{solution}

\begin{solution}
If $s,t$ agree on a basis cover $\{B_i\}$ of $U$, then they agree on an ordinary open cover of $U$. The sheaf locality axiom yields $s=t$.
\end{solution}

\begin{solution}
Compatible local solutions $u_i$ glue to a unique smooth $u$. On each $U_i$, $Lu=Lu_i=0$. Since this is local and the $U_i$ cover $U$, $Lu=0$ everywhere. Locality is inherited from the smooth-function sheaf.
\end{solution}